\chapter{Proyectos y planificación de producción}
\label{chap:planificacion}

\section{CPM: precedencias y camino crítico}
CPM trabaja sobre una red acíclica de actividades: una actividad solo empieza
cuando han terminado todas sus predecesoras. No asigna recursos limitados.
El recorrido hacia delante calcula inicios y fines tempranos; el recorrido
hacia atrás calcula los tardíos. La holgura total es inicio tardío menos
inicio temprano. Una actividad crítica tiene holgura cero.

En \path{manual/examples/plan_cpm.json}, A dura 3 días, B dura 2 y ambas
preceden a C, que dura 4. Cada actividad tiene costo normal 10, tiempo de
crash una unidad menor y costo de crash 15.
\begin{lstlisting}[language=bash]
swift run qsb validate-project-json manual/examples/plan_cpm.json
swift run qsb solve-project-json manual/examples/plan_cpm.json --backend native
swift run qsb solve-cpm-crash-json manual/examples/plan_cpm.json 6 --backend native
\end{lstlisting}
\begin{figure}[htbp]
\centering
\begin{tikzpicture}
\node[qsb activity,line width=1.4pt] (A) at (0,1) {A\\3};
\node[qsb activity] (B) at (0,-1) {B\\2};
\node[qsb activity,line width=1.4pt] (C) at (4,0) {C\\4};
\draw[qsb selected] (A) -- (C);
\draw[qsb alternative] (B) -- (C);
\end{tikzpicture}

\caption{Red CPM de actividades: el número bajo cada nombre es su duración en días. El trazo grueso identifica actividades y precedencias críticas; el discontinuo muestra la otra precedencia.}
\label{fig:cpm}
\end{figure}

El proyecto normal termina en 7 días. A y C son críticas; B tiene holgura 1.
El costo normal total es 30. Crash para plazo 6 reduce C de 4 a 3, con costo
incremental 5 y costo total 35. Reducir A también podría alcanzar el mismo
costo; no suponga unicidad del plan.

La optimización de crash supone costos lineales entre tiempos normal y mínimo,
con pendiente $(C_{crash}-C_{normal})/(t_{normal}-t_{crash})$ cuando hay reducción
posible. Usa tiempos continuos y un LP educativo que considera caminos;
el número de caminos puede crecer con rapidez. No es una política de
contratación de recursos ni incluye costo indirecto del proyecto.
Los plazos inferiores a lo alcanzable deben tratarse como infactibles, nunca
como instrucciones para violar precedencias.

\section{PERT: duración esperada e incertidumbre}
PERT utiliza estimaciones optimista $a$, más probable $m$ y pesimista $b$,
ordenadas $a\leq m\leq b$. La aproximación clásica calcula
\[
 E[T]=(a+4m+b)/6,\qquad \operatorname{Var}(T)=((b-a)/6)^2.
\]
En \path{plan_pert.json} las ternas son A(2,3,4), B(1,2,3) y C(3,4,5),
con las mismas precedencias del ejemplo CPM.
\begin{lstlisting}[language=bash]
swift run qsb solve-project-json manual/examples/plan_pert.json --backend native
\end{lstlisting}
\begin{figure}[htbp]
\centering
\begin{tikzpicture}
\node[qsb activity,line width=1.4pt] (A) at (0,1) {A\\(2,3,4)\\$E[T]=3$};
\node[qsb activity] (B) at (0,-1) {B\\(1,2,3)\\$E[T]=2$};
\node[qsb activity,line width=1.4pt] (C) at (4,0) {C\\(3,4,5)\\$E[T]=4$};
\draw[qsb selected] (A) -- (C);
\draw[qsb alternative] (B) -- (C);
\end{tikzpicture}

\caption{Red PERT: cada actividad muestra la terna optimista, más probable y pesimista, y su duración esperada en días. El trazo grueso identifica el camino crítico calculado con esas medias.}
\label{fig:pert}
\end{figure}

La programación de tiempos esperados vuelve a durar 7, con A y C críticas.
La varianza reportada es $2/9$ y la desviación estándar $0.471405$ días.
En este ejemplo hay un único camino crítico y la suma de sus varianzas
presupone independencia. La implementación suma las varianzas de todas las
actividades marcadas críticas; si hay varios caminos críticos, ese total no
es la varianza de un único camino ni la del máximo entre ellos. El camino crítico
puede cambiar entre realizaciones; sustituir duraciones aleatorias por sus
medias no calcula exactamente la esperanza del máximo. La etiqueta
\texttt{exact} describe el cálculo sobre esos tiempos esperados, no una
predicción probabilística exacta de la fecha final. QSB no devuelve una
probabilidad de cumplimiento ni simula esa incertidumbre en este módulo.

\subsection{CPM y PERT en macOS}
Importe el sobre de proyecto o cree una definición CPM/PERT. El editor de
actividades recoge nombres, predecesoras y tiempos; CPM añade costos y PERT
las tres estimaciones. Predecesoras inexistentes y ciclos aparecen en
Validation. Cambiar de método conserva los datos de ambos métodos después
de su inicialización: compruebe las estimaciones al volver a PERT.
Copy Table/Paste Table permite trasladar actividades desde una hoja de cálculo.
Run produce duración, actividades críticas, barras tempranas/tardías, holguras
y tabla temporal. Revise JSON para la varianza y los valores exactos.
La operación CLI de crash es un flujo separado; la vista normal no implica
que se haya optimizado un plazo. Native y Validate están disponibles;
External no resuelve proyectos actualmente.

\section{Planificación agregada}
La planificación agregada equilibra demanda por períodos mediante producción
regular, horas extra, subcontratación, inventario y pedidos pendientes.
Cuando la capacidad depende de plantilla, añade contratación y despido.
Un balance básico es
\[
 I_t-B_t=I_{t-1}-B_{t-1}+R_t+O_t+S_t-D_t.
\]
El objetivo suma costos de producción, capacidad ociosa, personal, inventario
y atraso según los datos. QSB normaliza los métodos históricos simple,
transporte y LP a una misma formulación LP continua; el campo \texttt{method}
conserva su identidad de entrada, no selecciona tres algoritmos nativos distintos.

El ejemplo \path{plan_aggregate.json} tiene dos meses, demanda 4 y 6,
capacidad regular 5 cada mes, costo regular 2 por unidad, horas extra hasta 2
con costo 5, inventario con costo 1 y sin atraso ni subcontratación.
No depende de plantilla y empieza sin inventario.
\begin{lstlisting}[language=bash]
swift run qsb validate-aggregate-json manual/examples/plan_aggregate.json
swift run qsb solve-aggregate-json manual/examples/plan_aggregate.json --backend native
\end{lstlisting}
El plan fabrica 5 en cada mes; conserva 1 al final del primero y cero al final
del segundo. Cuesta $10\cdot2+1=21$. Fabricar exactamente la demanda del primer
mes y usar una unidad extra en el segundo costaría $9\cdot2+5=23$.
Las unidades deben ser consistentes: capacidad por trabajador no es capacidad
total y \texttt{capacityRequirementPerUnit} convierte producción a capacidad.
Los vectores deben tener el mismo número de períodos; \texttt{null} en máximos
opcionales representa ausencia de cota. La plantilla y las cantidades son
continuas, incluso si el negocio necesita personas enteras. No redondee una
solución sin comprobar de nuevo balances y capacidad.

En macOS el editor de períodos separa demanda, costos, inventario/fuentes y
plantilla. Copy Table transfiere las 20 columnas de períodos; Paste Table valida
el rango antes de sustituirlo. Run muestra producción apilada y demanda,
inventario, atrasos y tabla del plan. Un atraso aparece bajo cero en el gráfico,
pero su cantidad del modelo es no negativa. Native resuelve el LP educativo;
External usa HiGHS para la misma formulación si está instalado. Validate no
produce un plan. El objetivo óptimo corresponde a los costos y cotas introducidos,
no a criterios ausentes como estabilidad social de la plantilla.

\section{MRP: explosión de necesidades}
MRP parte de un programa maestro de producción (MPS), una lista de materiales
(BOM), inventario, recepciones programadas y plazos. Recorre padres antes que
componentes: la liberación de una orden del padre genera necesidades brutas
para sus componentes, multiplicadas por el uso unitario \cite{silver1998inventory}.
Para un artículo con stock de seguridad $SS$, necesidad bruta $G_t$ y
recepción programada $SR_t$, la necesidad neta es
$\max(0,G_t+SS-I_{t-1}-SR_t)$. Una recepción planificada cubre esa necesidad,
y su liberación se desplaza hacia atrás por el plazo.

\path{plan_mrp.json} tiene semanas S0--S3. Cinco unidades de P deben
recibirse en S3; P utiliza dos C por unidad. Ambos tienen plazo de una semana,
stock inicial y seguridad cero, regla LFL y capacidad ilimitada.
\begin{lstlisting}[language=bash]
swift run qsb validate-mrp-json manual/examples/plan_mrp.json
swift run qsb solve-mrp-json manual/examples/plan_mrp.json --backend native
\end{lstlisting}
P recibe 5 en S3 y libera 5 en S2. C necesita 10 en S2, recibe 10 entonces y
libera 10 en S1. Los inventarios proyectados quedan en cero. Las necesidades
de C provienen de la liberación de P, no de la fecha de demanda final.
La BOM debe ser acíclica y todos sus identificadores deben existir.

Las reglas disponibles son LFL, EOQ, LUC, LTC y PPB. LFL recibe la necesidad
neta; EOQ parte de $\sqrt{2DS/H}$ y recibe al menos la necesidad actual.
En EOQ, \texttt{lotSizingLookahead} se interpreta como múltiplo de lote
\(m=\max(1,\text{valor configurado})\), no como número de períodos:
la cantidad es \(\max(\text{necesidad},m\lceil EOQ/m\rceil)\).
Si el costo de posesión anual es cero, recibe la necesidad neta.
LUC compara costo unitario; LTC y PPB usan actualmente el mismo selector,
que minimiza la diferencia absoluta entre costo de posesión y costo de pedido
dentro del horizonte configurado. Son reglas hacia delante, no una
optimización global de costos multinivel. QSB marca el documento MRP como
\texttt{heuristic}, incluso cuando el caso particular usa LFL.
La capacidad finita genera \texttt{capacityExcess}; no nivela ni reprograma
órdenes. Si la liberación cae antes del horizonte, QSB la acumula en el primer
bucket mediante $\max(0,t-L)$; ese registro requiere revisión operativa y no
prueba que el plazo pueda cumplirse.

En macOS Item Master recoge parámetros, Schedules cantidades por período y
Bill of Materials relaciones padre/componente. Seleccione el artículo para
editar o leer su calendario; los horizontes largos se agrupan para navegar.
Copy/Paste Table siempre intercambia el rango completo de su sección.
La solución muestra necesidades, inventario, recepciones, liberaciones y excesos
de capacidad por artículo, con JSON como auditoría. Native y Validate están
disponibles; External no está implementado para MRP.

\section{Ejercicios}
\begin{enumerate}
\item En CPM aumente B de 2 a 4. Recalcule camino crítico y holguras.
\item En PERT aumente solo el pesimista de A. Compare duración y varianza;
      explique qué supuestos limitan una interpretación probabilística.
\item Eleve el costo de inventario agregado de 1 a 4. Compare producir
      anticipadamente con usar horas extra.
\item Añada inventario inicial 4 de C en MRP y compruebe su necesidad neta.
      Después limite la capacidad de C en S2 a 5 y explique el exceso.
\end{enumerate}
